\scanpage{42}
% L1–2 | top, small note above the bracket
Catégorie des $k$-Algèbres $[$\uline{graduées}$]$
% slanted
\margin{\unsure{sera souvent|pas néc.} \unsure{sous-entendu|ss.-entendu}} bivariantes (\inserted{anticommutatives,} munies d'un degré $d(A)$, nulles en degré $\notin[1,d(A)]$ ; on ne donne pas de $k_A:A-k$).

% L3–5 | upper
Objet particulier de cette catégorie : $k$ de degré 0.

Algèbre graduée bivariante jaugée : munie d'un hom dans $k$. Dans certains cas d'hom de telles, $f_*$ est déterminé par $f^*$. $(*)$

% L6–8 | upper
Produit tensoriel d'Algèbres bivariantes (resp. jaugées). Dans certains cas, c'est un produit dans la catégorie des Algèbres bivariantes jaugées.

% L9–12 | upper middle, between and beneath lines
Algèbre graduée bivariante \inserted{(resp. gr. biv. jaugée)} \uline{simpliciale}. \deleted{\unsure{Exemp}} Cas \uline{spécial} : on a $A(I'\amalg I'')\xleftarrow[\sim]{(\alpha^*,\beta^*)}A(I')\otimes A(I'')$, d'autre part si $\alpha:I-J$, $\beta:I'\to J'$, alors $\alpha\amalg\beta=\unsure{\gamma|r}:I\amalg I'\longrightarrow J\amalg J'$, on veut que $(\alpha\amalg\beta)_*:A(\unsure{J\amalg J'|I\amalg I'})-A(I\amalg I')$ s'identifie à $\alpha_*\otimes\beta_*$.

% L12–17 | middle
On voit que si $A(\Delta^0)=A_0$ est ``de Poincaré'' (cas jaugé) $A(\deleted{\unsure{*}})-A(0)$ est une équiv. des $A(*)$ \unsure{sur spéc|sur spés} de Poincaré, et des alg. de Poincaré. Et $A(*)-(A(0),\delta)$ est un foncteur pl. fid de la catégorie des Alg. graduées bivariantes jaugées simpliciales, dans la catégorie des algèbres \inserted{gr.} bivariantes jaugées $A$, munies d'un $\Delta\in[A\otimes A]^{d(A)}$ : $[$il convient de préciser les axiomes à mettre sur ce $\Delta]$
% slanted
\margin{\uline{à vérifier} ; il faut d'abord noter qu'il suffit de connaître $\delta_*:A\longrightarrow A\otimes A$ ; on prend $\Delta=\delta_*(1)$}

% L18–25 | lower, example and left margin
\uuline{Ex} 1) Si on a une théorie ch. \deleted{sur la \unsure{variété}} $[$satisfaisant Künneth et Poincaré$]$ sur la catégorie $\unsure{\mathcal V|V}$\margin{$k$ un \uline{corps}} des variétés complètes \unsure{équidim.} non singulières de dim $d$\inserted{, sur un corps donné}\deleted{, \unsure{de dim}, \illegible{1}} d'où un foncteur \inserted{cov} de $\mathcal V$ dans la catégorie $C_{\mathrm{jaug}}$ des Alg. gr. bivar. jaugées sur $k$, et en transformant par ce foncteur l'objet simplicial canonique $I\longmapsto X^I$ défini par $X\in\mathcal V$, on trouve une Algèbre \deleted{bivariante} simpliciale \deleted{bivariante} jaugée canonique, d'ailleurs définie aussi canoniquement, comme dit plus haut, par \unsure{*|le} \unsure{uniquement|\deleted{correspondants}|composant} $H^*(X)$\ldots.

% L26–27 | footnote under horizontal rule
$(*)$ Attention, les \uline{isom} $(f^*,f_*)$ ne sont pas néc. tels que $f_*=f^{*-1}$ i.e. $f_*(1)=1$ ($\Rightarrow$ \deleted{\unsure{tels}} que $k_A f^*=k_B$), ex : $f^*=\mathrm{id}$, $f_*(y)=y\unsure{v|\upsilon|\sigma}$ $(v\in A^*)$, $k_B(y)=k_A(yv^{-1})$ \ldots
% archive sheet number: 41

